\chapter{Conclusion}
\label{ch:conclusion}

NeuroTensor developed a complete, versioned course pipeline from public
resting-state \EEG{} to tensor representation, low-rank decomposition,
participant-level classification, and artifact replay. The work began by solving
a practical data-engineering problem: \FTD{} limited both people and duration.
An \FTD{}-referenced MILP match produced a controlled 69-person cohort while
preserving all excess windows. Signal audit, CAR, deterministic outlier rules,
and boundary guards then made the analysis partitions explicit and reproducible.

Three tensor views were tested: a non-negative
subject-channel-frequency-time tensor, signed logPSD channel-frequency-local-
segment tensors, and raw-waveform channel-lag-position delay tensors summarized
by exact second-order moments. CP, graph NCP, Tucker/HOSVD, TT, and supervised TT
were evaluated alongside spectral features and standard classifiers.

The highest retained family result is 59.7\% mean balanced accuracy for tuned
RBF SVM, while the V4 primary complete-pipeline endpoint is 54.5\%. \CN{} is the
easiest class and \FTD{} remains the hardest. Tucker is competitive in the
corrected 69-person experiment, but spectral controls outperform waveform-only
methods in the 88-person experiment. Fusion is promising but does not isolate a
TT effect. XGBoost is useful yet not best; longer optimization improves neither
the fixed supervised TT nor the evidence base by itself.

The project's strongest contribution is therefore methodological: it shows how
to move from attractive tensor ideas to an auditable experiment while preserving
negative results and corrections. The current models are meaningful course
benchmarks, not clinical systems. The next credible advance is new independent
\FTD{} data and a locked external-validation protocol, not another unbounded
search on the same participants.

\begin{keybox}{Final takeaway}
Tensor decomposition provides a principled language for preserving multi-way
EEG structure, but trustworthy performance depends equally on cohort design,
subject-level validation, training-only transformations, optimization diagnostics,
and honest endpoint reporting.
\end{keybox}

